\documentclass[10pt]{article}
\usepackage[T1]{fontenc}
\usepackage[utf8]{inputenc}
\usepackage[french]{babel}
\usepackage{amsmath,amssymb,mathrsfs,mathpazo}
\usepackage{geometry}
\geometry{a4paper,margin=2cm}
\usepackage{graphicx,array,tabularx,tabularray,colortbl,multirow,multicol,booktabs,enumitem,pifont,ulem,fancybox,diagbox,cancel}
\usepackage{pstricks,pstricks-add,pst-plot,pst-tree,pst-node,pst-text,pst-func,pst-math,pst-bspline,pst-3dplot,pst-eucl,pst-solides3d}
\let\pscalculate\undefined
\usepackage{pst-bezier}
\usepackage{tikz,pgfplots,tkz-tab}
\usetikzlibrary{arrows,arrows.meta,patterns,calc,positioning,trees,decorations.markings,shapes.geometric,shapes.misc,shapes.arrows,intersections,decorations.pathmorphing,decorations.pathreplacing,backgrounds,shadows,babel}
\pgfplotsset{compat=1.18}
\usepgfplotslibrary{fillbetween}
\usepackage[np]{numprint}
\usepackage{listings,makecell,pas-tableur,twoopt}
\let\barre\undefined
\usepackage[pstricks]{tablvar}
\definecolor{darkgreen}{RGB}{0,100,0}
\providecommand{\red}{\color{red}}
\providecommand{\blue}{\color{blue}}
\providecommand{\green}{\color{green}}
\providecommand{\black}{\color{black}}
\providecommand{\white}{\color{white}}
\providecommand{\gray}{\color{gray}}
\providecommand{\cyan}{\color{cyan}}
\providecommand{\magenta}{\color{magenta}}
\providecommand{\yellow}{\color{yellow}}
\providecommand{\violet}{\color{violet}}
\providecommand{\orange}{\color{orange}}
\providecommand{\euro}{\texteuro}
\providecommand{\EUR}{\texteuro}
\providecommand{\ieme}{\textsuperscript{e}}
\providecommand{\degres}{\ensuremath{^\circ}}
\providecommand{\degree}{\ensuremath{^\circ}}
\providecommand{\vv}[1]{\overrightarrow{#1}}
\providecommand{\Coord}[1]{#1}
\providecommand{\decosix}{}
\providecommand{\starredbullet}{\ensuremath{\bullet}}
\pagestyle{empty}
\setlength{\parindent}{0pt}
\usepackage{adjustbox}\DeclareUnicodeCharacter{2236}{:}\DeclareUnicodeCharacter{2190}{\ensuremath{\leftarrow}}
\providecommand{\R}{\mathbb{R}}
\providecommand{\N}{\mathbb{N}}
\providecommand{\Z}{\mathbb{Z}}
\providecommand{\Q}{\mathbb{Q}}
\providecommand{\D}{\mathbb{D}}
\providecommand{\C}{\mathbb{C}}
\providecommand{\e}{{\mathrm{e}}}
\providecommand{\vect}[1]{\overrightarrow{#1}}
\providecommand{\vectt}[1]{\overrightarrow{\mathrm{#1}}}
\providecommand{\vvt}[1]{\overrightarrow{\mathrm{#1}}}
\providecommand{\np}[2][]{{#2~#1}}
\providecommand{\vv}[1]{{\overrightarrow{#1}}}
\renewcommand{\d}{{\mathrm{d}}}
\providecommand{\degres}{{^\circ}}
\providecommand{\degre}{{^\circ}}
\providecommand{\blue}{\color{blue}}
\providecommand{\red}{\color{red}}
\providecommand{\enskip}{\quad}
\providecommand{\llbracket}{\lbrack\!\lbrack}
\providecommand{\rrbracket}{\rbrack\!\rbrack}
\providecommand{\euro}{{\text{€}}}
\providecommand{\no}{{\text{n°}}}
\providecommand{\interval}[2]{{[#1;#2]}}
\providecommand{\intervFO}[2]{{[#1;#2[}}
\providecommand{\dotfill}{\ldots}

\providecommand{\ds}{\displaystyle}
\providecommand{\textding}[1]{#1}
\providecommand{\Oij}{\(\mathrm{O};\vec{\imath},\vec{\jmath}\)}
\providecommand{\Oijk}{\(\mathrm{O};\vec{\imath},\vec{\jmath},\vec{k}\)}
\providecommand{\Ouv}{\(\mathrm{O};\vec{u},\vec{v}\)}
\providecommand{\intervFF}[2]{\ensuremath{[#1;#2]}}
\renewcommand{\euro}{\ifmmode\text{\texteuro}\else\texteuro\fi}
\DeclareUnicodeCharacter{00B0}{\ifmmode^{\circ}\else\textdegree\fi}
\DeclareUnicodeCharacter{20AC}{\ifmmode\text{\texteuro}\else\texteuro\fi}
% Énoncés : Éducation nationale. Transcription et corrigés : APMEP, Denis Vergès et contributeurs. Découpage et classement : CoopMaths, Sébastien Lozano, Éric Elter et contributeurs. Adaptation de la présentation : Maths Explorer.
% CC BY-SA : https://coopmaths.fr/www/about/
% https://coopmaths.fr/alea/static/bac/2022/tex/bac_2022_05_sujet1_etrangers_1.tex
\begin{document}
\noindent\textit{APMEP / CoopMaths. CC BY-SA. Présentation adaptée par Maths Explorer.}\par\medskip
\medskip

\emph{Cet exercice est un questionnaire à choix multiples. Pour chacune des questions suivantes, une seule des quatre réponses proposées est exacte. Les six questions sont indépendantes.}

\smallskip

\emph{Une réponse incorrecte, une réponse multiple ou l'absence de réponse à une question ne rapporte ni n'enlève de point. Pour répondre, indiquer sur la copie le numéro de la question et la lettre de la réponse choisie. \\
Aucune justification n'est demandée.}

\medskip

\begin{enumerate}
\item On considère la fonction $f$ définie pour tout réel $x$ par $f(x) = \ln \left(1  + x^2\right)$. 

Sur $\R$, l'équation $f(x) = \np{2022}$

\begin{center}
\begin{tabularx}{\linewidth}{*{2}{X}}
\textbf{a.~~} n'admet aucune solution. &\textbf{b.~~} admet exactement une solution.\\
\textbf{c.~~}admet exactement deux solutions.&\textbf{d.~~} admet une infinité de solutions.
\end{tabularx}
\end{center}

\item Soit la fonction $g$ définie pour tout réel $x$ strictement positif par: 
\[g(x) = x \ln (x) - x^2\]

On note $\mathcal{C}_g$ sa courbe représentative dans un repère du plan.

\begin{center}
\begin{tabularx}{\linewidth}{*{2}{X}}
\textbf{a.~~} La fonction $g$ est convexe sur $]0~;~+\infty[$.&\textbf{b.~~} La fonction $g$ est concave sur $]0~;~+\infty[$.\\
\textbf{c.~~}La courbe $\mathcal{C}_g$ admet exactement un point d'inflexion sur $]0~;~+\infty[$.&\textbf{d.~~}La courbe $\mathcal{C}_g$ admet exactement
deux points d'inflexion sur $]0~;~+\infty[$.
\end{tabularx}
\end{center}

\item On considère la fonction $f$ définie sur $]- 1~;~1[$ par

\[f(x) = \dfrac{x}{1 - x^2}\]

Une primitive de la fonction $f$ est la fonction $g$ définie sur l'intervalle $] - 1~;~1[$ par :

\begin{center}
\begin{tabularx}{\linewidth}{*{2}{X}}
\textbf{a.~~} $g(x) = - \dfrac12 \ln \left(1 - x^2\right)$&\textbf{b.~~} $g(x) =  \dfrac{1 + x^2}{ \left(1 - x^2\right)^2}$\\
\textbf{c.~~} $g(x)= \dfrac{x^2}{2\left(x - \dfrac{x^3}{3}\right)}$&\textbf{d.~~} $g(x) = \dfrac{x^2}{2}\ln \left(1 - x^2\right)$
\end{tabularx}
\end{center}

\item La  fonction $x \longmapsto  \ln \left(-x^2- x + 6\right)$ est définie sur

\begin{center}
\begin{tabularx}{\linewidth}{*{2}{X}}
\textbf{a.~~} $]- 3~;~2[$&\textbf{b.~~} $]- \infty~;~6]$\\
\textbf{c.~~} $]0~;~+\infty[$&\textbf{d.~~} $]2~;~+\infty[$
\end{tabularx}
\end{center}

\item On considère la fonction $f$ définie sur $]0,5~;~+ \infty [$ par 

\[f(x) =x^2- 4x+ 3 \ln (2x - 1)\]

Une équation de la tangente à la courbe représentative de $f$ au point d'abscisse 1 est:

\begin{center}
\begin{tabularx}{\linewidth}{*{2}{X}}
\textbf{a.~~}$y = 4x- 7$ &\textbf{b.~~}$y = 2x - 4$\\
\textbf{c.~~}$y = -3(x - 1) + 4$ &\textbf{d.~~}$y = 2x - 1$
\end{tabularx}
\end{center}

\item L'ensemble $S$ des solutions dans $\R$ de l'inéquation $\ln (x + 3) < 2\ln (x + 1)$ est:

\begin{center}
\begin{tabularx}{\linewidth}{*{2}{X}}
\textbf{a.~~}$S = ]- \infty~;~-2[ \cup ]1~;~+\infty[$&\textbf{b.~~}$S = ]1~;~+ \infty[$\\
\textbf{c.~~} $S = \emptyset$&\textbf{d.~~} $S = ]- 1~;~1[$
\end{tabularx}
\end{center}

\end{enumerate}

\bigskip

\end{document}
